\section{Sviluppi futuri}

Lo sviluppo futuro più naturale è la costruzione di uno scenario più verosimile, con veicoli eterogenei, traffico misto, geometrie stradali più complesse, ostacoli e modelli di propagazione più ricchi. Un'estensione di questo tipo permetterebbe di capire meglio quanto la robustezza osservata dipenda dal meccanismo \emph{SafeSwitch} e quanto invece dalle semplificazioni dello scenario.

Un secondo sviluppo riguarda il confronto più controllato tra LTE e 5G, mantenendo fissi scenario, parametri radio e repliche statistiche. In questo modo sarebbe possibile isolare meglio l'effetto della migrazione tecnologica e misurare anche metriche aggiuntive, come latenza, jitter e durata degli handover.

In questo contesto, l’implementazione del 5G proposta rappresenta dunque una base solida e riutilizzabile per futuri lavori sperimentali di ricerca, abilitando lo studio di scenari basati sull’integrazione di reti eterogenee quali 802.11p, VLC e 5G.
